% Irodov, Part II "Thermodynamics and Molecular Physics",
% section 2.2 "The First Law of Thermodynamics. Heat Capacity";
% original numbered problems 26--35 (printed pp.65--66 / PDF pp.74--75).
% Gas molar masses M and heat capacities are ideal-gas approximations.

\thermoenergyproblem{26}{65}
证明：若室内空气与外界相通、外界压强
$p$ 保持恒定，则\emph{室内空气的总内能}
与温度无关。
若室内容积 $V=40\,\mathrm{m^3}$，
空气压强为标准大气压，
估计室内空气的内能 $U$。
\exsolution
将室内空气视为双原子理想气体，
暂不激发振动自由度，
其绝热指数 $\gamma=C_p/C_V\simeq7/5$。
室内气体总内能为
\[
U=\nu C_VT,
\]
其中 $\nu$ 是室内气体的摩尔数。
因为理想气体状态方程
$pV=\nu RT$，且
$C_V=R/(\gamma-1)$，
故
\[
\boxed{U=\frac{pV}{\gamma-1}}.
\]
当房间容积 $V$ 和外界压强 $p$ 固定，
$U$ 不随温度变化；升温时部分空气流出，
使摩尔数 $\nu\propto1/T$。
代入 $\gamma-1=0.4$、
$p=1.01325\times10^5\,\mathrm{Pa}$、
$V=40\,\mathrm{m^3}$：
\[
\boxed{U\simeq\frac{1.01325\times10^5\cdot40}{0.4}
\simeq1.0\times10^7\,\mathrm J=10\,\mathrm{MJ}}.
\]
\excheck 这一结论要求\emph{房间与外界通气}；
密闭、质量不变的房间内理想气体内能仍随温度变化。

\thermoenergyproblem{27}{66}
一个绝热封闭容器中装有理想气体，
其摩尔质量为 $M$、
热容比为 $\gamma=C_p/C_V$。
容器及气体原先整体以速率 $v$
平动，突然停止。
若停止过程的气体宏观平动动能
最终完全转化为其内能，求气体升温 $\Delta T$。
\exsolution
设气体总质量为 $m$，摩尔数
$\nu=m/M$。
停止前相对实验室的气体宏观平动动能为
$K=\frac12mv^2$。
绝热且假定没有把能量带到容器外，
这部分宏观动能转化为热内能：
\[
\nu C_V\Delta T=\frac12mv^2.
\]
由
$C_V=R/(\gamma-1)$、
$\nu=m/M$，消去质量：
\[
\frac mM\frac{R}{\gamma-1}\Delta T
=\frac12mv^2,
\]
得到
\[
\boxed{\Delta T=\frac{(\gamma-1)Mv^2}{2R}}.
\]
\excheck 这不是任意急停过程的无条件能量守恒结论；
假定气体原有宏观平动能全部内化，
且容器壁、外部制动机构不吸收相应份额。

\thermoenergyproblem{28}{66}
两只绝热刚性容器由带阀门的短管相连，
内部装的是同一种理想气体（空气）。
两容器容积、初始压强和温度分别为
$(V_1,p_1,T_1)$ 与 $(V_2,p_2,T_2)$。
打开阀门，最终达到共同平衡温度和压强。
求末温 $T$ 与末压 $p$。
\exsolution
设两容器最初摩尔数为
\[
\nu_i=\frac{p_iV_i}{RT_i}\quad(i=1,2).
\]
整个系统刚性、绝热，外界没有机械功
与热量交换，因此总内能守恒：
\[
(\nu_1+\nu_2)C_VT
=\nu_1C_VT_1+\nu_2C_VT_2.
\]
消去同一种气体的 $C_V$，
\[
T=\frac{\nu_1T_1+\nu_2T_2}{\nu_1+\nu_2}
=\boxed{\frac{p_1V_1+p_2V_2}
{p_1V_1/T_1+p_2V_2/T_2}}
=\boxed{\frac{T_1T_2(p_1V_1+p_2V_2)}
{p_1V_1T_2+p_2V_2T_1}}.
\]
总摩尔数守恒且总体积为 $V_1+V_2$，
于是
\[
p(V_1+V_2)=(\nu_1+\nu_2)RT
=p_1V_1+p_2V_2,
\]
即
\[
\boxed{p=\frac{p_1V_1+p_2V_2}{V_1+V_2}}.
\]
\excheck 平衡温度是初始\emph{摩尔数加权平均}，
不简单等于 $(T_1+T_2)/2$。

\thermoenergyproblem{29}{66}
密闭容器装有初始标准状况的氢气，
容积为 $V=5.0\,\mathrm L$，
将其冷却
$\Delta T=55\,\mathrm K$。
求气体内能变化以及气体放出的热量。
\exsolution
标准状况取
$T_0\simeq273.15\,\mathrm K$、
$p_0\simeq1.013\times10^5\,\mathrm{Pa}$。
理想氢气在此温度视为双原子气体，
$C_V=5R/2=R/(\gamma-1)$，
其中 $\gamma\simeq7/5$。
密闭刚性容器体积不变，
因此气体没有对外体积功。
初始摩尔数
$\nu=p_0V/(RT_0)$，
降温的内能变化
\[
\Delta U=\nu C_V(-\Delta T)
=-\frac{p_0V\Delta T}{T_0(\gamma-1)}.
\]
代入 $V=5.0\times10^{-3}\,\mathrm{m^3}$：
\[
\boxed{\Delta U\simeq-2.5\times10^2\,\mathrm J
=-0.25\,\mathrm{kJ}}.
\]
热力学第一定律
$Q=\Delta U+A$，等容过程 $A=0$，
故气体\emph{吸收}的热为 $Q=\Delta U<0$，
向环境放出的热量大小为
\[
\boxed{Q_{\rm out}=-Q=-\Delta U\simeq0.25\,\mathrm{kJ}}.
\]
\excheck 区分带符号的气体吸热 $Q$ 与
放热量的正值 $Q_{\rm out}$。

\thermoenergyproblem{30}{66}
对氮气进行等压加热，
使气体在过程中对外做功
$A=2.0\,\mathrm J$。
假设氮气为常温双原子理想气体，
求需提供的热量 $Q$。
\exsolution
等压理想气体体积功
$A=\nu R\Delta T$。
吸热为
$Q=\nu C_p\Delta T$。
相除：
\[
\frac QA=\frac{C_p}{R}.
\]
由 $C_p-C_V=R$ 和
$\gamma=C_p/C_V$，
得 $C_p=\gamma R/(\gamma-1)$。
因此
\[
\boxed{Q=\frac{\gamma}{\gamma-1}A}.
\]
氮气 $\gamma\simeq7/5$，
故
\[
\boxed{Q=\frac{1.4}{0.4}(2.0\,\mathrm J)
\simeq7.0\,\mathrm J}.
\]
\excheck 供热不仅成为外功，还使气体内能增加；
所以 $Q>A$。

\thermoenergyproblem{31}{66}
一摩尔某理想气体经等压加热，
温度上升 $\Delta T=72\,\mathrm K$，
吸热 $Q=1.60\,\mathrm{kJ}$。
求气体做功 $A$、内能增量 $\Delta U$
和热容比 $\gamma=C_p/C_V$。
\exsolution
等压理想气体做功
\[
\boxed{A=R\Delta T
=(8.314)(72)\,\mathrm J\simeq0.60\,\mathrm{kJ}}.
\]
由第一定律
$Q=\Delta U+A$，
\[
\boxed{\Delta U=Q-A\simeq1.00\,\mathrm{kJ}}.
\]
定压摩尔热容
$C_p=Q/\Delta T$，
定容摩尔热容
$C_V=(Q-R\Delta T)/\Delta T$，
故
\[
\boxed{\gamma=\frac{C_p}{C_V}
=\frac{Q}{Q-R\Delta T}
\simeq1.60}.
\]
\excheck 用 $\gamma\simeq1.60$ 反算
$C_V=R/(\gamma-1)\simeq13.9\,\mathrm{J/(mol\,K)}$，
与 $\Delta U/\Delta T$ 一致。

\thermoenergyproblem{32}{66}
某理想气体共 $\nu=2.0$ mol，
初温 $T_0=300\,\mathrm K$。
先等容冷却，使压强降低
$n=2.0$ 倍；
再等压膨胀，使温度恢复到初始值 $T_0$。
求整个过程从外界\emph{净吸收}的热量。
\exsolution
等容第一段由 $p\propto T$，
末温为 $T_1=T_0/n$。
第二段等压加热使温度
从 $T_1$ 回到 $T_0$。
两段吸热总和为
\[
Q=\nu C_V(T_1-T_0)
+\nu C_p(T_0-T_1).
\]
由 $C_p-C_V=R$，得
\[
Q=\nu R(T_0-T_1)
=\nu RT_0\left(1-\frac1n\right).
\]
因此
\[
\boxed{Q=(2.0)(8.314)(300)(1-1/2)
\simeq2.5\,\mathrm{kJ}}.
\]
\excheck 第一段放热、第二段吸热，
\emph{净吸热}为正，且不依赖气体热容比。

\thermoenergyproblem{33}{66}
某理想混合气由
$\nu_1=2.0$ mol 氧气、
$\nu_2=3.0$ mol 二氧化碳组成。
求混合气的热容比 $\gamma$，
以两气体各自热容比
$\gamma_1,\gamma_2$ 表示并给出估算值。
\exsolution
混合气的总定容热容为各组分热容之和：
\[
C_{V,\Sigma}=
\nu_1\frac R{\gamma_1-1}
+\nu_2\frac R{\gamma_2-1}.
\]
总定压热容同样相加：
\[
C_{p,\Sigma}=
\nu_1\frac{\gamma_1 R}{\gamma_1-1}
+\nu_2\frac{\gamma_2 R}{\gamma_2-1}.
\]
所以混合气的热容比为
\[
\boxed{\gamma=
\frac{\nu_1\gamma_1(\gamma_2-1)
+\nu_2\gamma_2(\gamma_1-1)}
{\nu_1(\gamma_2-1)+\nu_2(\gamma_1-1)}}.
\]
如果使用原书温区下的常见近似
$\gamma_1\simeq1.40$（氧气）、
$\gamma_2\simeq1.30$（二氧化碳），
\[
\boxed{\gamma\simeq
\frac{2(1.40)(0.30)+3(1.30)(0.40)}
{2(0.30)+3(0.40)}
=\frac{2.40}{1.80}\simeq1.33}.
\]
\excheck CO$_2$ 的有效热容比依赖温度和
是否激发振动自由度；题目未给具体温度，
数值结果需说明所采用的组分热容假设。

\thermoenergyproblem{34}{66}
由质量为 $7.0\,\mathrm g$ 的氮气和
$20\,\mathrm g$ 的氩气组成理想混合气。
在常温、热容近似不随温度变化时，
求混合气的质量比热容
$c_V,c_p$。
\exsolution
氮气为双原子气体，摩尔热容约
$C_{V,{\rm N_2}}=5R/2$、
$C_{p,{\rm N_2}}=7R/2$；
氩气为单原子气体，分别为
$3R/2$ 与 $5R/2$。
两种气体的摩尔数：
\[
\nu_{\rm N_2}=\frac{7.0}{28}=0.25\,\mathrm{mol},
\qquad
\nu_{\rm Ar}=\frac{20}{40}=0.50\,\mathrm{mol}.
\]
混合气总质量 $m_{\Sigma}=27.0\,\mathrm g$。
以每克质量计的定容比热为
\[
c_V=
\frac{\nu_{\rm N_2}(5R/2)+\nu_{\rm Ar}(3R/2)}
{m_{\Sigma}}
=\frac{(0.625+0.75)R}{27.0\,\mathrm g},
\]
故
\[
\boxed{c_V\simeq0.42\,\mathrm{J/(g\,K)}}.
\]
同样
\[
c_p=
\frac{\nu_{\rm N_2}(7R/2)+\nu_{\rm Ar}(5R/2)}
{m_{\Sigma}}
=\frac{(0.875+1.25)R}{27.0\,\mathrm g},
\]
所以
\[
\boxed{c_p\simeq0.655\,\mathrm{J/(g\,K)}}.
\]
\excheck 组分按\emph{摩尔数乘摩尔热容}求总热容，
最后除以总质量；不能直接对质量比热做无权平均。

\thermoenergyproblem{35}{66}
一摩尔理想气体在竖直圆筒中，
被无质量、无摩擦活塞封住，
活塞上方连通大气。
初温为 $T$，
现缓慢向上提起活塞，
使气体\emph{等温}膨胀到原体积的
$n$ 倍。
求外界提起活塞需要做的功 $A_{\rm ext}$。
\exsolution
气体初始处于平衡，
因此初始气压
$p_{\rm gas}(V_0)=p_0$。
一摩尔理想气体
$p_{\rm gas}(V)=RT/V$，
且初态 $p_0V_0=RT$。
缓慢提拉活塞的外力大小为
$F_{\rm ext}=[p_0-p_{\rm gas}(V)]S$，
其中 $S$ 为活塞横截面积，
允许外力由向上逐渐变大。
若活塞上升 $h$，体积变化
$\mathrm dV=S\,\mathrm dh$，
因此外力所做功
\[
A_{\rm ext}
=\int_{V_0}^{nV_0}
\left(p_0-\frac{RT}{V}\right)\mathrm dV
=p_0V_0(n-1)-RT\ln n.
\]
由 $p_0V_0=RT$，
\[
\boxed{A_{\rm ext}=RT(n-1-\ln n)}.
\]
\excheck 对 $n\ge1$，
不等式 $n-1\ge\ln n$ 保证外力功非负；
它不同于气体自身的等温膨胀功 $RT\ln n$。
